\section{Related Work}
\label{sec:related}

The $k$-regret query~\cite{NS10}, also known as the regret-minimizing set query, addresses the limitations of top-$k$ and skyline queries. It returns $k$ representative tuples such that, for any unknown nonnegative linear utility function over $d$ attributes, the best returned tuple closely approximates the global optimum. Its objective is to minimize the maximum regret ratio defined in Section~\ref{sec:def}, thereby avoiding explicit utility specification while controlling the output size. However, finding an optimal $k$-regret set is NP-hard~\cite{AK17, CL17, CT14}, particularly in high dimensions. Early methods, including \textsc{Cube}~\cite{NS10}, \textsc{Greedy}~\cite{NS10}, and \textsc{GeoGreedy}~\cite{PW14}, either lack theoretical guarantees or perform poorly in practice. Later approaches such as \textsc{HD-Greedy}~\cite{asudeh2017efficient} and $\varepsilon$-\textsc{Kernel}~\cite{CL17} improve efficiency but offer weak bounds: \textsc{HD-Greedy} provides only a loose upper bound, while the bound of $\varepsilon$-\textsc{Kernel} may exceed $1$ even for $d<10$, making it uninformative in practice.

Recent advances have extended regret minimization to high-dimensional settings. Xie et al.~\cite{XW18} proposed the \textsc{Sphere} algorithm, which achieves an asymptotically optimal upper bound on the maximum regret ratio without restricting data dimensionality. Although earlier methods such as Cube ~\cite{NS10} also provide theoretical guarantees that apply across dimensions, Sphere is distinguished by offering an asymptotically optimal bound while also demonstrating strong empirical performance, making it a state-of-the-art single-round solution. However, its performance degrades as dimensionality increases due to two limitations: (1) it requires the output size $k$ to be at least the dimensionality $d$, which is impractical in high-dimensional applications, and (2) its runtime grows sub-exponentially with $d$, rendering it inefficient for large datasets. These issues motivate the development of complementary approaches for scalable high-dimensional decision making.

Another line of work studies interactive regret minimization, which replaces one-shot query processing with multi-round interaction to progressively infer the user’s implicit utility function. The first such framework, \textsc{UtilityApprox}~\cite{NL12}, repeatedly presents a small set of tuples and asks the user to select the preferred one. Each response refines the estimated utility function, allowing subsequent queries to focus on more promising tuples and reduce regret for a fixed output size~\cite{NL12}. However, \textsc{UtilityApprox} relies on synthetic tuples, which may frustrate users when a preferred option is not real, and requires an impractically large number of questions in high-dimensional settings.

Xie et al.~\cite{XCL19} addressed this limitation with a strongly truthful framework that displays only real tuples, improving user trust and simplifying preference elicitation. Their algorithms, including \textsc{UH-Random} and \textsc{UH-Simplex}, converge to near-optimal results within a bounded number of rounds and achieve an asymptotically optimal question complexity in two dimensions. However, their scalability is limited in higher dimensions by numerical instability in convex-hull computations. To overcome this issue, these methods require users to explicitly specify a small subset of relevant attributes, which imposes substantial burden in high-dimensional settings, as will be justified in later paragraphs.

Several other interactive methods have also been proposed. \textsc{Active-Ranking}~\cite{jamieson2011active} adaptively selects pairwise comparisons to recover a complete ranking under a Euclidean embedding assumption. However, it targets full-ranking recovery rather than regret minimization and assumes that the data are in general position, which may not hold in practical databases. \textsc{Preference-Learning}~\cite{qian2015learning} similarly reconstructs the entire utility function through adaptive comparisons, requiring queries that may be irrelevant to identifying the user’s preferred tuples. More recently, Wang et al.~\cite{WWX21} proposed \textsc{HD-PI} and \textsc{RH} for identifying top-$k$ tuples. Although these methods provide dimension-independent guarantees under ideal conditions, their scalability deteriorates as dimensionality increases because convex-hull and skyline frontiers grow rapidly, increasing both computational and interaction costs. Chen and Wong~\cite{CW23} further study best-tuple retrieval under random interaction errors through \textsc{Verify-Point} and \textsc{Verify-Space}. These methods provide probabilistic retrieval guarantees over a predefined attribute space, whereas FHDR discovers key attributes implicitly from high-dimensional data.

A further limitation of existing works is upfront attribute selection. For example, a user may be asked to select $5$ attributes from $100$ directly, before seeing actual candidate tuples. Do et al.~\cite{LA25} simplify interaction by comparing two attributes at a time, but still require users to select their attributes of interest. Prior work often treats this step as trivial preprocessing, while our survey and user study indicate that upfront attribute selection can be burdensome.

This selection can impose a substantial cognitive burden: users must inspect and compare many attributes while working memory is limited~\cite{Cowan_2001}. Moreover, users’ preferences are often constructed during actual decision-making rather than specified in advance~\cite{Slovic1995}, so selection without concrete candidate tuples risks prematurely excluding relevant attributes. On the other hand, partial-profile designs, which show only a subset of attributes at a time, have been reported to reduce decision complexity and choice randomness while improving response efficiency~\cite{meyerhoff_performance_2023}. Comparisons between concrete tuples also offer a more intuitive way to elicit preferences~\cite{NL12, XCL19, WWX21, CW23}. For example, a user may be unable to state abstractly whether ``location'' is globally more important than ``floor area'', yet can readily choose between two houses once the trade-off is instantiated. In this way, comparison-based interaction reveals preferences implicitly with less cognitive effort. These observations motivate FHDR’s implicit discovery of key attributes through tuple comparisons.

Compared with existing works, FHDR natively supports high-dimensional regret minimization and maintains strong performance as dimensionality increases, addressing key scalability barriers. Its user-friendly design requires no upfront attribute selection and allows users to skip or stop answering questions at any time. Table~\ref{tab:capabilities} summarizes the advantages of our method, alongside comparison with existing methods. By combining high-dimensional scalability, computational efficiency, and flexible interaction, FHDR advances the state of the art and enables practical and accurate decision support for complex, real-world multi-criteria applications.

\begin{table*}[!t]
  \centering
  \small
  \setlength{\tabcolsep}{3pt}
  \renewcommand{\arraystretch}{1.15}
  \begin{tabular}{@{}p{0.23\textwidth}*{5}{>{\centering\arraybackslash}p{0.14\textwidth}}@{}}
    \toprule
    \textbf{Method} & \textbf{Practical high-dimensional support} & \textbf{Exact-favorite guarantee} & \textbf{Implicit key-attribute discovery} & \textbf{User-controlled output size} & \textbf{Output at any interaction budget} \\
    \midrule
    \textsc{Sphere}~\cite{XW18} & \halfcheckmark & $\times$ & $\times$ & $\times$ & $\checkmark$ \\
    \textsc{UtilityApprox}~\cite{NL12} & \halfcheckmark & $\checkmark$ & $\times$ & $\times$ & $\times$ \\
    \textsc{UH-Random}~\cite{XCL19} & $\times$ & $\checkmark$ & $\times$ & $\times$ & $\times$ \\
    \textsc{HD-PI} / \textsc{RH}~\cite{WWX21} & $\times^{\dagger}$ & $\checkmark$ & $\times$ & $\checkmark$ & $\times$ \\
    \textsc{Verify-Point}~\cite{CW23} & $\times^{\dagger}$ & $\checkmark$ & $\times$ & $\checkmark$ & $\times$ \\
    \textsc{Breakpoint}~\cite{LA25} & \halfcheckmark & $\times$ & $\times$ & $\times$ & $\checkmark$ \\
    \midrule
    \textbf{FHDR} & $\checkmark$ & $\checkmark$ & $\checkmark$ & $\checkmark$ & $\checkmark$ \\
    \bottomrule
  \end{tabular}
  \caption{Capabilities of representative methods. Practical high-dimensional support reflects practical computation and interaction costs in the evaluated high-dimensional settings; \halfcheckmark\ denotes partial support with substantial practical limitations. $\dagger$ Execution entries concern the tested representatives \textsc{HD-PI} and \textsc{Verify-Point}, respectively. Detailed qualifications and reproduction results appear in \iftoggle{TechnicalReport}{Section~\ref{appen:capabilities} of the \technicalreport}{Appendix~\ref{appen:capabilities}}.}
  \label{tab:capabilities}
\end{table*}
